\scanpage{67}
% loc: 1–3, top
et on trouve que
\[
 C^*=\coprod_{e\in\underbrace{\underline{\omega}^*/k^{*2}}_{\substack{\text{torseur}\\\text{sous }k^\times/k^{\times2}}}}C^*(e)
\]

% loc: 4–11
partition qui est permutée transitivement par les homothéties \inserted{de $M$} : \deleted{l'on \unsure{que|pour}} $\lambda C^*(e)=C^*(\lambda e)$. L'ens des classes de cette partition \deleted{est} \deleted{\unsure{prend donc un}} de façon naturelle (par la relation d'homothétie) un torseur sous $\unsure{k^{\times2}/k^{\times2}|k^\times/k^{\times2}}$, et ce torseur est can. iso à $\underline{\omega}^\times/k^{\times2}$.

% loc: 12–16
Cette sous-division de $C^\times$ en classes ne dépend que de $M$, et pas du choix de $E$ (avec $\det E\simeq k$) et d'un iso $\operatorname{End}(E)\simeq M$. On peut aller plus loin, \deleted{\unsure{quand on On}}.

% loc: left at 12–17, with arrow to Non
% slanted
\margin{On doit supposer ici $ {}_2\operatorname{Pic}(k)=0$ ; (Que se passe-t-il autrement ?) Il semble que \unsure{même en|Mais on a} $=0$ \unsure{de|pour} cette subdivision de $C^*$ est invariante par le groupe \unsure{privé|propre} $\operatorname{Aut}(M)$ \unsure{des|de}. $\longrightarrow$ \uuline{Non}}

% loc: 17–19
On \deleted{\unsure{aura}} a une application surjective
\[
 \begin{aligned}\varphi^e:\ E^*&\longrightarrow C^*(e)\\
 x&\longmapsto\varphi^e(x\otimes x)\end{aligned}
\]

% loc: 20–22, bottom
satisfaisant \unsure{à}
\[
 \left\{\begin{aligned}
 \varphi^e(\lambda x)&=\lambda^2\varphi^e(x)\\
 \varphi^{\lambda e}(x)&=\lambda\varphi^e(x)
 \end{aligned}\right.
\]
% archive sheet number: 66

